% Extended Data figure: horizon analysis (generated by horizon_study/make_ed_horizon.py; values in
% horizon_study/tables_no_v4pro/ed_horizon_values.md; slopes from horizon_study/tables_no_v4pro/pooled.md, section 4).
\begin{figure}[htbp]
\centering
\includegraphics[width=\linewidth]{Figures/ed_horizon.png}
\caption{\textbf{Post-cutoff performance by forecasting horizon.}
The horizon is the number of months from the knowledge cutoff of a model to the first-observable date of an event.
The figure shows only post-cutoff events (horizon of one month or more) of the full corpus (January 2024 to August 2026) in the base condition.
We put the events into five horizon bins (1--3, 4--6, 7--12, 13--24 and 25 or more months).
Each point is the mean of one bin, at the mean horizon of the bin.
Bands are 95\% confidence intervals from 10,000 event-level bootstrap samples.
We do not show bins with fewer than 20 items.
\textbf{a}, Feasibility Assessment: balanced accuracy, the mean of the accuracies on original and perturbed questions; dashed line, chance (0.5).
\textbf{b}, Mechanistic Forecasting: MCQ accuracy; dashed line, chance (0.25).
\textbf{c}, Temporal Prediction: absolute error of the predicted month.
\textbf{d}, Solution Design: FRQ rubric score; dashed line, pass threshold (5).
We fitted a mixed model with random intercepts for event and model.
In this model, MCQ accuracy decreases by 0.17 log-odds for each year of horizon (95\% CI 0.05 to 0.29; $p = 0.004$, Holm-corrected $p = 0.06$).
The FRQ score does not change with the horizon ($+$0.01 rubric points for each year, 95\% CI $-$0.01 to $+$0.04).
The date error decreases with the horizon, but this decrease is an artefact.
The models predict dates long after their cutoff, and a constant prediction gives the same decrease.}
\label{fig:ed_horizon}
\end{figure}
